\chapter{Which targets are learnable?}\label{ch:learnability}

Part V turns from describing the data to deciding what to build on it. Before a model is built, these questions must be answered: is there a signal at all, out of time and not merely in sample; is the signal more than a single feature or a
simple rule already captures; and does the answer hold across markets and across an independent replicate of the data?
The country series answered them for $\est{ch18.targets}{value}$ candidate targets in $\est{ch18.scopes}{value}$ scopes. This chapter examines that learnability map,
fits the transparent model the series recommended for its one learnable target, and identifies a field whose values carry
no information by construction. \Cref{ch:readiness-gates} then turns the answers into gates.

\section{The protocol}\label{sec:le-protocol}

Each target is predicted from features computed strictly before a cutoff (\cref{eq:mf-pit}) and judged on later cutoffs only.
The series reports the ROC AUC with a Hanley--McNeil interval \autocite{hanley1982} and the average precision as a multiple of
the base rate, the lift, and turns them into a verdict: no evidence of signal when the interval's lower bound does not exceed chance level; a weak signal, for rules rather than models, when the AUC is below $\est{ch18.rule-weak-auc}{value}$ or the lift
below $\est{ch18.rule-weak-lift}{value}$; a moderate signal below $\est{ch18.rule-learnable-auc}{value}$; learnable above. Controls give the verdicts weight. The fraud score shipped with the data, which \cref{sec:ad-fraud} showed to encode the
label, is entered as a target's feature: a protocol that finds no signal anywhere might simply be unable to find any, and this
positive control shows it can. And every target is evaluated again on the backup copy, an independent replicate
(\cref{ch:statistical-equivalence}), so that a verdict that holds in both is not an accident of one draw.

One caution applies to the intervals. A customer appears at several cutoffs, and the Hanley--McNeil interval treats every row
as independent, so the intervals are somewhat too narrow; the verdicts near a threshold should be read with that in mind.

\section{The map}\label{sec:le-map}

\Cref{tab:ch18-map} shows the out-of-time AUC of the $\est{ch18.targets}{value}$ targets in the $\est{ch18.scopes}{value}$
scopes. Of the $\est{ch18.cells}{value}$ evaluations, $\est{ch18.no-signal}{value}$ show no evidence of signal, and
$\est{ch18.verdict-stable}{value}$ of the targets receive the same verdict in every scope. The positive control works: with the
leaked score, fraud is predicted with AUCs between $\est{ch18.auc-min-leak}{value}$ and $\est{ch18.auc-max-leak}{value}$.
Without it, fraud lies between $\est{ch18.auc-min-fraud}{value}$ and $\est{ch18.auc-max-fraud}{value}$, and complaints within the following months ($\est{ch18.auc-min-complaint}{value}$ to $\est{ch18.auc-max-complaint}{value}$), delinquency
($\est{ch18.auc-min-dpd}{value}$ to $\est{ch18.auc-max-dpd}{value}$) and escalation after a contact
($\est{ch18.auc-min-escalation}{value}$ to $\est{ch18.auc-max-escalation}{value}$) show no signal anywhere. Campaign conversion
is moderate ($\est{ch18.auc-min-conversion}{value}$ to $\est{ch18.auc-max-conversion}{value}$, lift $\est{ch18.ap-lift-conversion}{value}$),
and the model adds a median of $\est{ch18.gain-conversion}{value}$ in AUC over its best single feature, the channel-by-segment
structure of \cref{sec:pc-campaigns}: a contact rule, not a model.

\input{../build/generated/tables/ch18-map}

Dormancy is the one learnable target: AUCs between $\est{ch18.auc-min-dormant}{value}$ and $\est{ch18.auc-max-dormant}{value}$,
with a median lift of $\est{ch18.ap-lift-dormant}{value}$. But the number of transactions in the previous $\est{ch18.months-past}{value}$ months alone
reaches between $\est{ch18.single-min-dormant}{value}$ and $\est{ch18.single-max-dormant}{value}$. The boosted model learns a
customer's activity rate, the exposure of \cref{sec:ag2-exposure}, and nothing else.

\section{Dormancy as a personal Poisson rate}\label{sec:le-poisson}

If a customer transacts as a Poisson process with a personal monthly rate $\lambda$, and the rates vary across customers as a
gamma distribution with shape $a$ and rate $b$, the prediction of dormancy has a closed form.

\begin{proposition}\label{prop:le-gp}
Let $\lambda \sim \mathrm{Gamma}(a, b)$, let $k \mid \lambda \sim \mathrm{Poisson}(s \lambda)$ be the count over the $s$ months
before a cutoff, and let the next $u$ months be independent of the past given $\lambda$. Then
\begin{equation}\label{eq:le-gp}
  \Pr(\text{no transaction in the next } u \text{ months} \mid k) = \Bigl(\frac{b + s}{b + s + u}\Bigr)^{a + k}.
\end{equation}
\end{proposition}

\begin{proof}
By conjugacy, $\lambda \mid k \sim \mathrm{Gamma}(a + k, b + s)$. The probability of no event in $u$ months given $\lambda$ is
$e^{-u\lambda}$, and its expectation under that posterior is the gamma's Laplace transform at $u$,
$\bigl((b + s)/(b + s + u)\bigr)^{a + k}$.
\end{proof}

The marginal of $k$ is the negative binomial of \textcite{greenwood1920}, so $a$ and $b$ are fitted by maximum likelihood on
the training cutoffs, truncated to the active customers the target concerns ($k \ge 1$), with $s = \est{ch18.months-past}{value}$
and $u = \est{ch18.months-future}{value}$. Rebuilt here from the transactions, the target has $\est{ch18.rows-train}{value}$
training and $\est{ch18.rows-test}{value}$ test rows, and the fit gives $a = \est{ch18.gp-a}{value}$ and
$b = \est{ch18.gp-b}{value}$, a mean rate of $\est{ch18.gp-mean-rate}{value}$ transactions a month.

The model with its pair of parameters ranks the test customers with an AUC of $\est{ch18.gp-auc}{estimate}$ ($\est{ch18.gp-auc}{ci}$),
indistinguishable from the boosted model's $\est{ch18.nb-dormant-auc}{value}$ ($\est{ch18.nb-dormant-lo}{value}$ to
$\est{ch18.nb-dormant-hi}{value}$). It is also calibrated, which a ranking metric cannot show (\cref{tab:ch18-calibration}): its
mean prediction is $\est{ch18.gp-mean-pred}{value}$ against an observed $\est{ch18.gp-obs}{value}$, and no decile is off by more
than $\est{ch18.cal-gap-max}{value}$. Its Brier score \autocite{brier1950} is $\est{ch18.gp-brier}{value}$ against
$\est{ch18.gp-brier-base}{value}$ for predicting the base rate, a skill of $\est{ch18.gp-brier-skill}{value}\,\%$. The shrinkage
of the gamma prior is what does the work: the naive plug-in $e^{-u k / s}$, which treats each customer's observed rate as known,
predicts a mean of $\est{ch18.naive-mean-pred}{value}$ and scores worse than the base rate (skill
$\est{ch18.naive-brier-skill}{value}\,\%$). The champion for this target is a model with a pair of parameters that explains every prediction
in a sentence: a customer with $k$ recent transactions has the dormancy probability of \cref{eq:le-gp}.

\input{../build/generated/tables/ch18-calibration}

\section{A field drawn at random}\label{sec:le-random}

A learnability map can only be as good as the fields its features read, and one field of the complaints is empty of information
by construction. Of the $\est{ch18.cp-total}{value}$ complaints, $\est{ch18.cp-with-product}{value}$ name an affected product,
and every one of them, $\est{ch18.cp-other-owner}{value}$, names a product of another customer. If the product were drawn
uniformly from the bank's products, the share of complaints naming a product of the complainant's own market would equal that
market's share of all products. It does, to within $\est{ch18.cp-max-gap}{value}$ in every market, and a goodness-of-fit test
does not reject it ($p = \est{ch18.cp-chi2-p}{value}$). The affected product is a random draw; the integrity rule that flags it
(\cref{sec:ag-policy}) measures the generator, not the bank, and any feature built from it is noise.

\section{Verdict and implications}\label{sec:le-verdict}

Of the candidate targets, one is learnable, one is a contact rule, and the rest show no signal in any market or in either copy,
while the positive control shows that the protocol would have found a signal. The learnable target is the arithmetic of personal
rates, and a gamma-Poisson model with a pair of parameters, calibrated and explainable, ranks it as well as boosting. This is the factual basis
of the readiness gates: they must certify not only that a target is learnable but that the learning beats the simplest model of
the same thing.

\section{Limitations}\label{sec:le-limits}

The learnability map is the country series' stored result, with intervals that ignore the repetition of customers across cutoffs.
The dormancy target is rebuilt here on processing months with quarterly cutoffs, close to but not identical with the series'
construction; the agreement of the AUCs suggests the difference is immaterial.
